\chapter*{List of symbols}
\addcontentsline{toc}{chapter}{List of symbols}

Symbols are listed in the order in which they first appear. Where the same letter is
used in the literature for two different things, the table says which meaning is used here.

\subsection*{Earthquakes (Chapter 1)}
\begin{longtable}{@{}p{22mm}p{120mm}@{}}
\toprule
Symbol & Meaning \\
\midrule
$\mu_{\mathrm s}$ & shear modulus (rigidity) of rock, about $3\times10^{10}$\,Pa \\
$A$ & area of the fault surface that slipped in an earthquake (m$^2$) \\
$\bar D$ & average slip, i.e.\ relative displacement of the two sides of the fault (m) \\
$M_0$ & seismic moment, $M_0=\mu_{\mathrm s} A\bar D$ (N\,m); the physical size of an earthquake \\
$M_w$ & moment magnitude, $M_w=\tfrac23(\log_{10}M_0-9.1)$; a dimensionless logarithmic size \\
$N(\geq M_w)$ & number of earthquakes (per year, in a region) with magnitude at least $M_w$ \\
$a$ & Gutenberg--Richter $a$-value: $\log_{10}$ of the number of events with $M_w\ge 0$; measures overall activity \\
$b$ & Gutenberg--Richter $b$-value: slope of $\log_{10}N$ against $M_w$; measures how fast the number of events falls with size \\
$\tau$ & exponent of a power-law probability density, $P(x)\propto x^{-\tau}$ \\
\bottomrule
\end{longtable}

\subsection*{Probability and critical phenomena (Chapter 2)}
\begin{longtable}{@{}p{22mm}p{120mm}@{}}
\toprule
Symbol & Meaning \\
\midrule
$P(x)$ & probability density (or, for integer $x$, probability) of a random quantity $x$ \\
$\bar F(x)$ & complementary cumulative distribution, the probability that the quantity is $\ge x$ \\
$\avg{x^q}$ & $q$-th moment, the average of $x^q$ \\
$x_{\min}$, $x_c$ & lower and upper cutoffs of a power law \\
$p$, $p_c$ & occupation probability in percolation and its critical value \\
$n_s$ & number of percolation clusters of $s$ sites, per lattice site \\
$\xi$ & correlation length: the largest distance over which the system is correlated \\
$\nu$, $\gamma$ & critical exponents of $\xi$ and of the mean cluster size \\
$D$ & fractal dimension: the size of an object of linear extent $R$ grows as $R^{D}$ \\
\bottomrule
\end{longtable}

\subsection*{Sandpile models (Chapter 2)}
\begin{longtable}{@{}p{22mm}p{120mm}@{}}
\toprule
Symbol & Meaning \\
\midrule
$L$ & linear size of the lattice (number of sites along one side) \\
$N$ & number of sites, $N=L^d$ in $d$ dimensions \\
$k$ & number of neighbours of a site in the bulk ($k=2$ on a chain, $4$ on the square lattice) \\
$z_i$ & height (number of grains) at site $i$ \\
$z_c$ & critical height; a site with $z_i\ge z_c$ is unstable. In this thesis $z_c=k$ \\
$s$ & avalanche size: total number of topplings caused by one added grain \\
$a$ & avalanche area: number of distinct sites that toppled (not to be confused with the GR $a$-value) \\
$T$ & avalanche duration: number of parallel update steps \\
$\Delta$ & toppling matrix ($N\times N$) \\
$G$ & Green function, $G=\Delta^{-1}$ \\
$n_{ij}$ & number of topplings at site $i$ caused by a grain added at site $j$ \\
$g_j$ & $\sum_i G_{ij}$, the mean avalanche size for a grain added at $j$ \\
$l$, $l'$ & position of the empty site in the 1D chain before and after an avalanche ($0$ means ``no empty site'') \\
$\sigma(q)$ & moment exponent, defined by $\avg{s^q}\propto L^{\sigma(q)}$ \\
\bottomrule
\end{longtable}

\subsection*{The Burridge--Knopoff model (Chapter 4)}
\begin{longtable}{@{}p{22mm}p{120mm}@{}}
\toprule
Symbol & Meaning \\
\midrule
$m$ & mass of a block \\
$k_c$, $k_p$ & spring constants of the coupling springs (between blocks) and of the pulling springs (to the plate) \\
$F_0$ & maximum static friction force \\
$v$, $v_1$ & velocity of the plate; velocity over which kinetic friction weakens \\
$X_j$, $U_j$ & displacement of block $j$; the same in units of $D_0$ \\
$\omega_p$, $\tau$ & $\sqrt{k_p/m}$; dimensionless time $\omega_pt$ \\
$D_0$ & $F_0/k_p$, the stretch of a pulling spring at the friction threshold \\
$\ell$ & stiffness ratio, $\ell^2=k_c/k_p$ \\
$\nu$ & dimensionless loading speed $v/(\omega_pD_0)$ \\
$\alpha$ & rate of velocity weakening, $2\alpha=\omega_pD_0/v_1$ \\
$\sigma$ & drop of friction at the onset of slip (1991 law) \\
$\phi(y)$ & dimensionless friction as a function of the sliding velocity $y$ in units of $v_1$ \\
$a$ & spacing between neighbouring blocks (not the avalanche area of Chapter 2) \\
$\xi$ & stiffness length $\ell a$; also the sound speed in units of length per unit $\tau$ \\
$\omega_1$ & $\sqrt{2\ell^2+1}$, frequency of one block slipping in a stuck chain \\
$\epsilon$ & fraction of the released force lost into the pulling spring, $1/(2\ell^2+1)$ \\
$\tilde\xi$, $\tilde M$ & localization length and the corresponding crossover moment \\
$M$, $\mu$ & moment of an event, $a\sum_j\delta U_j$, and its magnitude $\ln M$ \\
$\mathcal R(\mu)$, $\mathcal D(\mu)$ & number of events per unit length and per unit time (per unit plate displacement) \\
$b$, $b'$ & slopes of $\ln\mathcal D$ against $\mu$ for localized and delocalized events \\
$M_1$, $M^*$ & moment of a one-block event; upper cutoff of the delocalized events \\
\bottomrule
\end{longtable}
